\documentclass{article}
\usepackage[T1]{fontenc}
\usepackage{isabelle,isabellesym}
\usepackage{amsmath}
\usepackage{amssymb}
\usepackage[only,bigsqcap]{stmaryrd}

% this should be the last package used
\usepackage{pdfsetup}

% urls in roman style, theory text in math-similar italics
\urlstyle{rm}
\isabellestyle{it}

\begin{document}

\title{New Algebra: groups, rings, modules and fields with carrier sets}
\author{Lawrence C. Paulson}
\date{}

\maketitle

\begin{abstract}
  A development of abstract algebra in which every structure is a \emph{carrier
  set} together with its operations, following the locale-based style of
  Ballarin's \emph{Basic Algebra} \cite{ballarin-jacobson}, itself after
  Jacobson \cite{jacobson-basic-algebra}.  A group is a set with a composition
  and a unit; a ring is a set with two operations; an extension is a carrier with
  its own addition and multiplication.  Nothing is a record, and no structure is
  identified with a type, so substructures, quotients and extensions are all
  ordinary values that can be quantified over, ordered, and exhausted by
  transfinite arguments.

  The group theory covers subgroup lattices, normal series, the Zassenhaus and
  Schreier refinement theorems, the Jordan--H\"older uniqueness theorem, finite
  and indexed products, group actions, and solvability.  The ring theory covers
  ideals, quotients, the Chinese remainder theorem for finite families, Krull's
  theorem on the existence of maximal ideals, principal ideal domains, Euclidean
  domains, unique factorisation, and polynomials in one and in indexed families
  of indeterminates.  The module theory covers submodules, exact sequences,
  complements, free modules and their universal property, and finite-dimensional
  vector spaces with the Steinitz exchange lemma and a well-defined dimension.

  On these rest the field theory and Galois theory: minimal polynomials, the
  degree of a finite extension and the tower law, transitivity of algebraicity,
  normality and separability, Artin's theorem, the Galois correspondence, normal
  closures, the primitive element theorem, and the finite fields.  The
  cyclicity of the multiplicative group of a finite field is proved, along with
  the residue rings $\mathbb{Z}/n\mathbb{Z}$, their unit groups, and Euler's and
  Fermat's theorems.

  The development is self-contained: it depends only on the Isabelle/HOL
  distribution, and is intended to supersede the record-based \texttt{HOL-Algebra}
  rather than to build upon it.  Importing the theory \texttt{Algebra} makes the
  whole of it available.
\end{abstract}

% sane default for proof documents
\parindent 0pt\parskip 0.5ex

% for proof reading
\pagestyle{plain}

\tableofcontents
\newpage

% generated text of all theories
\input{session}

% optional bibliography
\bibliographystyle{abbrv}
\bibliography{root}

\end{document}
